\uuid{t0jE}
\exo7id{1823}
\titre{exo7 1823}
\auteur{drutu}
\organisation{exo7}
\datecreate{2003-10-01}
\isIndication{false}
\isCorrection{false}
\chapitre{Extremum, extremum lié}
\sousChapitre{Extremum, extremum lié}
\cpge{mpsi-mp2i.fonctions-deux-variables.b, pcsi.fonctions-deux-variables.b, ptsi.fonctions-deux-variables.b, mpsi-mp2i.fonctions-deux-variables.d, pcsi.fonctions-deux-variables.d, ptsi.fonctions-deux-variables.d, bcpst1.fonctions-deux-variables.c, mp-mpi.calcul-differentiel-optimisation.f}
\module{Calcul différentiel}
\niveau{L3}
\difficulte{}

\contenu{
\question{
Soit $f$ la fonction d\'efinie sur $\R^2$ par $f(x,y)=x^2-xy^2$.
Montrer que $(0,0)$ est le seul  point critique de $f$,  qu'il
n'est pas un extremum local, mais que pourtant la restriction de
$f$ \`a toute droite passant par $(0,0)$ admet en ce point un
minimum local.
}
}
